\section{Chain-of-Thought: convertir los Token de respuesta en trayectorias de razonamiento}

La segunda mitad de las diapositivas pasa de «qué puede hacer el modelo» a «cómo volver utilizables las capacidades que ya tiene». Un few-shot prompt estándar solo muestra question--answer; un CoT prompt muestra question--intermediate reasoning--answer, de modo que el modelo produce una serie de token intermedios componibles antes de generar la respuesta final. Es a la vez interface design y un aumento de la sequential inference computation.

\subsection{Motivación de la investigación: a una tarea compleja no se le puede dar solo medio segundo}

Una tarea sencilla puede requerir apenas un label, pero una task compleja de arithmetic, symbolic reasoning o commonsense necesita conservar estados intermedios. La intuición de la clase es esta: si una persona escribe los pasos al resolver un problema, exigir a un autoregressive model que emita de inmediato un final token equivale a no darle tiempo para desplegar el cálculo.

\lecturefigure{slide-20.jpg}{El artículo de Chain-of-Thought Prompting y el contexto de su demostración pública}{Diapositivas oficiales, p. 20; Wei et al., NeurIPS 2022}

Esta diapositiva de transición reúne el artículo, el equipo y la presentación en Google I/O, lo que muestra que CoT ya había pasado de prototipo de investigación a la atención del público. Para estas notas, lo importante no es el escenario del anuncio, sino el interface change, muy pequeño, de la siguiente figura: si se incluye o no la reasoning path en el exemplar.

\lecturefigure{slide-21.jpg}{Diferencia estructural entre standard prompting y chain-of-thought prompting}{Diapositivas oficiales, p. 21; Wei et al., 2022}

\begin{knowledgebox}{Lectura de la figura: lo único nuevo es un proceso intermedio que puede imitarse}
A la izquierda, el standard prompt solo muestra la question y la final answer, y ante un problema nuevo el modelo adivina directamente el resultado; a la derecha, el CoT prompt escribe los pasos de cálculo en la respuesta de ejemplo, así que el modelo continúa generando una reasoning path similar y después da la final answer. No se actualizan los weights ni se llama a una calculator externa; el cambio ocurre en el context y en la longitud de la generación.
\end{knowledgebox}

Desde la descomposición de probabilidad, la generación estándar de la respuesta puede escribirse como $p(a\mid q,e)$, donde $q$ es la pregunta y $e$ son los exemplars. CoT introduce una secuencia de token intermedios $z=(z_1,\ldots,z_m)$:
\[
p(a,z\mid q,e)=\prod_{i=1}^{m}p(z_i\mid q,e,z_{<i})\;p(a\mid q,e,z).
\]
La secuencia intermedia permite que el modelo escriba paso a paso en el context las multiplicaciones, las actualizaciones de estado y las subconclusiones; se trata de una intuición sobre el mecanismo, que no garantiza que los pasos textuales correspondan fielmente al cálculo interno ni que cada paso sea correcto.

\begin{lstlisting}[caption={Diferencia mínima entre un few-shot prompt estándar y un CoT prompt}]
standard = """Q: There are 15 trees and later 21. How many added?
A: 6
Q: 23 apples, use 20, buy 6. How many remain?
A:"""

cot = """Q: There are 15 trees and later 21. How many added?
A: 21 - 15 = 6. The answer is 6.
Q: 23 apples, use 20, buy 6. How many remain?
A:"""
\end{lstlisting}

\teachervoice{El ponente compara CoT con un profesor que pide a sus alumnos ``show your work''. Más adelante añade que no es justo darle a una persona un problema difícil y exigirle la respuesta en medio segundo; los token intermedios le dan al modelo espacio para desplegar un cálculo secuencial, pero por qué los pasos en lenguaje natural resultan tan eficaces seguía siendo, en 2023, una pregunta abierta.}

\subsection{Demo en vivo: el mismo problema pasa de 33 a 9}

La diapositiva ``CoT demo'' del deck oficial es solo un título, por lo que hay que volver a la grabación. En vivo se usa la misma arithmetic task: una cafetería tiene 23 apples, usa 20 y compra 6 más. Sin CoT exemplars, text-davinci-001 responde directamente 33; al cambiar a un preset con derivación intermedia, el modelo escribe $23-20=3$, $3+6=9$ y llega a la respuesta correcta.

\lecturefigure{slide-22.jpg}{La clase pasa de las diapositivas a la CoT demo en OpenAI Playground}{Diapositivas oficiales, p. 22}

Esta diapositiva de transición, casi vacía, no contiene resultados, pero marca un cambio en la fuente de la evidencia: las dos figuras siguientes no son ilustraciones del artículo, sino el live system state de la misma clase. La comparación debe conservar model, temperature, task y prompt preset; no basta con una sola captura del caso exitoso.

\videofigure{demo-no-cot-math-wrong.jpg}{Sin CoT: el modelo responde directamente 33 y no combina correctamente «usa 20, compra 6»}{Video oficial, 00:39:54}

\begin{knowledgebox}{Lectura de la figura: el fallo no es simplemente una suma mal hecha}
Los tres exemplar de arriba son question--short answer, y ante el problema nuevo el modelo mantiene el patrón de «dar directamente un número». El 33 probablemente proviene de combinar de forma errónea el 23 con parte de las variaciones; como no hay estados intermedios visibles, no podemos ubicar si el error ocurrió en la subtraction, en la addition o en la interpretación del enunciado. Este es justamente el punto ciego de diagnóstico de la only-final-answer interface.
\end{knowledgebox}

\videofigure{demo-cot-math-correct.jpg}{Con CoT: el modelo calcula primero que quedan 3, luego suma las 6 compradas y obtiene 9}{Video oficial, 00:40:17}

\begin{knowledgebox}{Lectura de la figura: el acierto proviene de actualizaciones de estado verificables}
Los exemplars del CoT preset escriben primero cada cambio de cantidad como una operación. La salida del problema nuevo conserva, en orden, la cantidad inicial, la cantidad consumida, el resto y lo añadido, así que el lector puede verificar los dos pasos locales $23-20=3$ y $3+6=9$. La figura respalda que «esta vez el prompt hizo que este modelo acertara y mostrara sus pasos», pero no demuestra que los pasos reflejen fielmente el proceso causal interno del modelo.
\end{knowledgebox}

La clase también muestra la last-letter concatenation: con el standard prompt, para ``Elon Musk'' el modelo produce `sk`; con CoT, primero toma la última letra de Elon, `n`, y la de Musk, `k`, y luego produce `nk`. Esto indica que el efecto no se limita a la arithmetic, sino que se extiende a cualquier task que pueda descomponerse en estados intermedios generables.

\subsection{Por qué la ganancia de CoT también emerge con la escala}

Tener pasos intermedios no garantiza que un modelo pequeño se beneficie. Un modelo pequeño puede acumular errores en una salida más larga, o simplemente no lograr abstraer la decomposition rule a partir del exemplar; un modelo grande tiene más probabilidad de generar correctamente cada paso, por lo que el delta entre CoT y standard prompting crece con la escala.

\lecturefigure{slide-23.jpg}{GSM8K y StrategyQA: la ganancia de CoT aparece en modelos suficientemente grandes}{Diapositivas oficiales, p. 23; Wei et al., 2022}

\begin{knowledgebox}{Lectura de la figura: comparar standard y CoT a la misma escala}
El eje horizontal es la escala del modelo; el vertical es, respectivamente, el performance en grade-school math y en commonsense reasoning. En los modelos pequeños ambos prompt dan resultados bajos, e incluso CoT no lleva ventaja; al entrar en la región de modelos más grandes, la curva de CoT sube con claridad y supera el fine-tuned state of the art de entonces. La evidencia central es el interaction effect que crece con la escala, no solo una puntuación alta aislada del modelo más grande.
\end{knowledgebox}

\begin{warningbox}{El texto de CoT no es automáticamente una explicación confiable}
El modelo puede llegar a la respuesta correcta y escribir pasos erróneos, o escribir pasos fluidos y terminar en una respuesta incorrecta. CoT mejora la observabilidad y el desempeño en la tarea, pero no puede tomarse directamente como mechanistic interpretability. Hace falta una verificación adicional con step verification, answer checking, counterfactual prompts o tools externas.
\end{warningbox}

\subsection{Resumen de la sección}

La esencia de CoT es que un frozen model genere token intermedios componibles antes de responder. La demo en vivo mostró la comparación sobre el mismo problema, y las curvas de GSM8K/StrategyQA mostraron la scale interaction. Es una inference interface, no una weight update; aumenta la sequential computation, pero no garantiza que la reasoning trace sea fiel ni correcta.
